\section{Related Work}

\subsection{Memory Mechanisms in LLMs}

\textbf{Context Window Extension.} Large language models (LLMs) are constrained by fixed-length context windows. Prior work extends context via sparse attention \cite{beltagy2020longformer,zaheer2020bigbird}, recurrence \cite{dai2019transformer,bulatov2022recurrent}, and length extrapolation \cite{chen2024clex,chen2025longpo}. However, longer context does not guarantee effective utilization: the ``Lost-in-the-Middle'' phenomenon persists \cite{liu2024lost, Bulatov2023ScalingTT}, suggesting context extension alone is insufficient for durable memory.


\textbf{Retrieval-Augmented and Parametric Memory.} Retrieval-augmented generation (RAG) \cite{lewis2020retrieval} externalizes memory to alleviate window limits, but its reliability depends on retrieval quality \cite{ram2023incontext}. Parametric approaches internalize information, yet often suffer from forgetting and instability \cite{delange2022continual}. Hybrid approaches \cite{wang2023longmem,packer2024memgpt} alleviate issues but lack a unified organizational principle for persistent memory.

\subsection{Memory Systems}
\textbf{Early Computational Memory.} Early differentiable memory systems (e.g., NTM/DNC/Key--Value memories) \cite{graves2014ntm,graves2016dnc,miller2016kv} introduced external memory interaction, but scale poorly and are ill-suited to modern autoregressive LLMs.

\textbf{Memory in LLM Agents.} As LLM-based agents evolve \cite{xi2023rise, xia2024agent}, memory systems have shifted toward persistent state integration. Recent systems introduce episodic \cite{wang2025mirix}, semantic \cite{shinn2024reflexion}, and hierarchical task memory \cite{sun2025hierarchical}. However, many designs still rely on fragmented text units and limited consolidation, which can degrade long-horizon performance \cite{packer2024memgpt}.

\textbf{Memory Operating Systems.} Recent work formalizes memory management as a system-level runtime. Some focus on lifecycle and capacity, such as \textit{Nemori}'s \cite{nemori2025} prediction-driven updates and \textit{MemoryOS}'s \cite{kang2025memory} hierarchical control. Others, like \textit{Mem0} \cite{mem02025} and \textit{Zep} \cite{zep2025}, prioritize structured fact maintenance via knowledge graphs, while \textit{MemOS} \cite{li2025memos} targets unified scheduling across memory types.

While these systems advance structural organization, they primarily focus on \textit{storage optimization} or \textit{fact maintenance}. EverMemOS distinguishes itself by implementing a \textit{three-phase memory lifecycle} that transforms episodic traces into synthesized semantic structures for long-horizon reasoning.
